% HISTORICAL DOCUMENT (checkpoint snapshot, pre-2026-06-10 evening).
% Superseded where it conflicts with GPU-MPC/ringlpn/CLAUDE.md and
% reports/dealerless_orca_ringlpn_full_proposal_2026_06_10.tex. In particular:
% real-OLE wiring and dense (slot) packing are NOW DONE; remaining oracle
% boundaries are SPFSS keygen and conversion wiring. Do not treat as current.

\documentclass[11pt]{article}

\usepackage[margin=1in]{geometry}
\usepackage{amsmath,amssymb,amsthm}
\usepackage{booktabs}
\usepackage{enumitem}
\usepackage{hyperref}
\usepackage{xcolor}

\hypersetup{
  colorlinks=true,
  linkcolor=blue!60!black,
  citecolor=blue!60!black,
  urlcolor=blue!60!black
}

\newcommand{\Z}{\mathbb{Z}}
\newcommand{\R}{\mathcal{R}}
\newcommand{\Party}{\mathsf{P}}
\newcommand{\ServerZero}{\mathsf{P}_0}
\newcommand{\ServerOne}{\mathsf{P}_1}
\newcommand{\Share}[1]{\langle #1 \rangle}
\newcommand{\modtwo}{\Z_{2^{\mathsf{bw}}}}
\newcommand{\PrimeRing}{\Z_p[X]/(X^N+1)}
\newcommand{\LimbRing}[1]{\Z_{p_{#1}}[X]/(X^N+1)}
\newcommand{\CRTRing}{\LimbRing{0} \times \LimbRing{1}}
\newcommand{\Orca}{\textsc{Orca}}
\newcommand{\RingLPN}{\textsc{Ring-LPN}}
\newcommand{\OLE}{\mathsf{OLE}}
\newcommand{\SPFSS}{\mathsf{SPFSS}}
\newcommand{\CRT}{\mathsf{CRT}}

\title{Toward a Dealerless \Orca{} Linear-Layer Protocol from \RingLPN{} OLE}
\author{GPU-MPC Ring-LPN Integration Notes}
\date{May 2026}

\begin{document}
\maketitle

\begin{abstract}
This note lays out a theory and engineering plan for replacing the trusted
\Orca{} fully-connected linear-layer dealer with a two-party preprocessing
protocol based on \RingLPN{}-style pseudorandom correlation generators
(PCGs) for oblivious linear evaluation (OLE)~\cite{boyle2022ringlpn}. The
current implementation already demonstrates an \Orca{}-compatible key format:
party key buffers are written in the existing \texttt{A}, \texttt{B},
\texttt{C\_masked} order and the online phase calls the unchanged
\texttt{gpuMatmulBeaver} routine. However, that checkpoint is still a
dealer/oracle correctness artifact. A scientifically sound dealerless result
should mean distributed seed/setup plus party-local expansion, not a hidden
central keywriter. It requires two additional protocol boundaries: (i) the
parties must obtain shares of the cross terms in a Beaver product without any
process seeing both parties' masks, and (ii) shares produced over a prime or
CRT modulus must be converted securely into \(\modtwo\), the arithmetic ring
consumed by \Orca{}. This document identifies the exact algebra, the current
validated evidence, the missing secure components, and a staged
implementation plan suitable for discussion before committing to the next
coding phase.
\end{abstract}

\section{Goal and Scope}

The near-term goal is not to rewrite \Orca{}'s online execution. The goal is
to replace the offline dealer work for linear layers while preserving the
online contract:
\[
  \texttt{gpuMatmulBeaver}
  \quad \text{consumes shares of} \quad
  (A, B, C_{\mathrm{masked}})
\]
where
\[
  C_{\mathrm{masked}} = A B + Y
  \quad \text{in} \quad \modtwo .
\]
Here \(A\) masks the input activation, \(B\) masks the weight matrix, and
\(Y\) is the output mask. The key files must remain byte-compatible with the
existing \Orca{} reader before we attempt model-level benchmarks.

The intended dealerless milestone is:
\[
  \ServerZero \text{ outputs } (A_0,B_0,C_0), \qquad
  \ServerOne \text{ outputs } (A_1,B_1,C_1),
\]
such that
\[
  A = A_0 + A_1,\qquad
  B = B_0 + B_1,\qquad
  C_0 + C_1 = AB + Y
  \pmod {2^{\mathsf{bw}}},
\]
and no trusted process learns both \(A_0,A_1\) or both \(B_0,B_1\).

\paragraph{Meaning of ``dealerless.''}
In the PCG literature, the trusted correlated-randomness source is replaced by
a setup protocol that distributes short correlated seeds, followed by local
``silent'' expansion of those seeds into long correlations
\cite{boyle2022ringlpn,boyle2019pcg}. Thus, ``dealerless'' should not be read
as zero setup or zero interaction. The intended claim is: after a secure
distributed setup, each party can locally expand its own seed/share and write
its \Orca{} key material, without any central process learning both parties'
linear masks or Beaver-product shares.

\paragraph{Out of scope for the first dealerless checkpoint.}
Convolution, ReLU/FSS, truncation, optimizer keys, and malicious-security
checks should remain separate unless the claim explicitly expands beyond
fully-connected linear layers. The first credible claim should be
\emph{dealerless FC matmul preprocessing}, not \emph{dealerless \Orca{} end to
end}.

\section{Current State in the Repository}

The current GPU-MPC workspace contains the following relevant artifacts.

\begin{table}[h]
\centering
\small
\begin{tabular}{p{0.29\linewidth}p{0.62\linewidth}}
\toprule
Artifact & Present status \\
\midrule
\texttt{bench\_ntt\_cuda} &
Promoted Cheddar-derived NTT/PolyMul backend supports requested
\(q=32,64,128\). Requested \(q=128\) is represented as two q62 CRT limbs and
reported as actual qbits \(124\). \\
\texttt{bench\_vole\_ringlpn} &
Standalone VOLE-style algebraic expansion prototype is validated, including
the q128 CRT path. This is an expansion benchmark, not full SPFSS-backed
dealerless key generation. \\
\texttt{bench\_ole\_ringlpn\_cuda} &
Standalone Figure 2 OLE artifact validates
\(z_0+z_1=x_0x_1\) over \(\Z_p[X]/(X^N+1)\). The source now has q128
argument plumbing via two q62 limbs, but saved summaries are primarily
q62/q64; q128 OLE summaries should be regenerated before making a
paper-parameter OLE claim. \\
\texttt{bench\_linear\_ole\_ringlpn\_cuda} &
Converts two OLEs into one ring-polynomial Beaver product and validates small
matrix semantics with shared operands. As with OLE, q128 should be rerun and
saved explicitly before claiming paper-parameter linear-layer evidence. \\
\texttt{test\_orca\_zp\_bridge} &
Host bridge validates carry-corrected export from \(\Z_p\), and from q128
CRT scalar modulus \(M=p_0p_1\), into \(\modtwo\) under explicit no-wrap
bounds. The conversion is dealer/oracle, not secure distributed conversion. \\
\texttt{bench\_orca\_fc\_ringlpn\_demo} &
Small \Orca{} FC demo writes byte-compatible \texttt{A}, \texttt{B},
\texttt{C\_masked} buffers. The current saved result includes q62 and q128
bounded cases and checks forward, \(dW\), and \(dX\) contracts. \\
\bottomrule
\end{tabular}
\caption{Relevant current implementation layers.}
\end{table}

The key scientific boundary is that the current \Orca{}-compatible demo is
still an oracle-style keywriter. It proves the \emph{format} and online
compatibility, but not yet the dealerless protocol.

\section{Scientific Soundness Audit}

The research-backed interpretation is summarized in Table~\ref{tab:audit}.
This table is meant to keep the presentation honest: the current codebase has
several strong correctness checkpoints, but the dealerless protocol claim
requires additional cryptographic functionality.

\begin{table}[h]
\centering
\small
\begin{tabular}{p{0.30\linewidth}p{0.30\linewidth}p{0.30\linewidth}}
\toprule
Question & Sound answer & Missing before a dealerless claim \\
\midrule
Can PCGs replace a trusted correlation dealer? &
Yes, in the PCG model: a distributed setup gives parties short correlated
seeds, then local expansion generates long OLE/triple correlations
\cite{boyle2022ringlpn,boyle2019pcg}. &
An actual setup/seed-generation boundary in the implementation, rather than a
central oracle that writes both parties' shares. \\
Does OLE fit Beaver triples? &
Yes. Two OLE-style product-share calls produce the cross terms
\(A_0B_1\) and \(A_1B_0\), completing the Beaver product algebra. &
A real OLE-backed transcript test that writes byte-compatible \Orca{} keys
from party-local state. \\
Is q128 a native Ring-LPN modulus here? &
No. In this codebase q128 means two q62 prime-limb computations with CRT
interpretation, not one monolithic prime-field Ring-LPN instance. &
Saved q128 OLE/linear summaries and a parameter note that states the exact
prime limbs and assumption variant. \\
Is the \(\Z_M \rightarrow \modtwo\) export local? &
No. The wrap/carry term depends on both additive shares. Modular conversion is
a separate MPC functionality~\cite{yu2011modcnv,escudero2020edabits}. &
A secure conversion protocol, or a redesign that avoids conversion before
\Orca{} consumes the key. \\
Can we cite generic Ring-LPN security? &
Not without parameters. NTT-friendly choices can make \(X^N+1\) reducible or
split, changing the assumption being invoked. &
A factorization/security audit for each \(p_i,N\), including whether the claim
uses irreducible, reducible, or splittable Ring-LPN. \\
\bottomrule
\end{tabular}
\caption{Scientific fit of the current plan against the relevant protocol
literature.}
\label{tab:audit}
\end{table}

\section{Algebraic Target}

Let \(A=A_0+A_1\) and \(B=B_0+B_1\) be additive shares of the input and weight
masks. A Beaver product expands as
\[
  AB =
  A_0B_0 + A_0B_1 + A_1B_0 + A_1B_1 .
\]
The local terms are easy:
\[
  \ServerZero \text{ can compute } A_0B_0,\qquad
  \ServerOne \text{ can compute } A_1B_1.
\]
The dealerless problem is exactly the cross terms:
\[
  A_0B_1 \quad \text{and} \quad A_1B_0.
\]
An OLE instance can be used to obtain additive shares of a product without
revealing the underlying operands; equivalently, symmetric OLE is commonly
viewed as secret-shared multiplication~\cite{baum2020ole}. We use the additive
convention \(U_0+U_1=A_0B_1\), noting that some papers use a sign convention
such as \(c-d=ab\). For each ring product, the protocol needs two OLE calls:
\[
\begin{aligned}
  (U_0,U_1) &\leftarrow \OLE(A_0,B_1), &
  U_0+U_1 &= A_0B_1,\\
  (V_0,V_1) &\leftarrow \OLE(A_1,B_0), &
  V_0+V_1 &= A_1B_0.
\end{aligned}
\]
Then the parties form
\[
\begin{aligned}
  Z_0 &= A_0B_0 + U_0 + V_0,\\
  Z_1 &= A_1B_1 + U_1 + V_1,
\end{aligned}
\]
which satisfies
\[
  Z_0+Z_1 = (A_0+A_1)(B_0+B_1).
\]
Finally, output-mask shares \(Y_0,Y_1\) are added:
\[
  C_0 = Z_0 + Y_0,\qquad C_1 = Z_1 + Y_1.
\]
These \(C_i\) are the \texttt{C\_masked} key entries consumed by
\texttt{gpuMatmulBeaver}.

\section{Where \RingLPN{} Enters}

The \RingLPN{} artifacts work over prime-limb negacyclic polynomial rings
\[
  \PrimeRing .
\]
For requested q128 in this codebase, the implementation should be interpreted
as two independent q62 prime limbs,
\[
  \CRTRing,
\]
with scalar CRT recomposition used only at the export boundary. This is more
precise than saying that the Ring-LPN assumption is directly over the
composite ring \(\Z_M[X]/(X^N+1)\), where \(M=p_0p_1\). The PCG literature
usually states assumptions and OLE constructions over prime fields/rings and
then uses splitting/CRT structure to obtain many scalar correlations
\cite{boyle2022ringlpn}.
This is structurally consistent with the concrete Ring-LPN PCG cost estimates
that use a modulus represented as a product of two 62-bit primes, but the
implementation still needs to state the exact prime-limb assumption and
security estimate rather than relying on that resemblance alone.

For the current correctness-first bridge, scalar values are embedded as
constant polynomials:
\[
  s \mapsto s + 0X + \cdots + 0X^{N-1}.
\]
This makes the first coefficient of the ring product equal to the scalar
product as long as the integer dot product does not wrap the modulus. For an
inner dimension \(K\) and bounded operands \(|x|,|w|\leq B\), the conservative
condition is
\[
  K B^2 < M .
\]
For q62, this is enough for small bounded demos but not for unrestricted
32-bit products. For q128/actual q124, the current host bridge validates the
bounded full-32-bit \(2\times2\times2\) case, but larger layers still need
layer-specific bounds.

\paragraph{Parameter assumption caveat.}
The current CUDA path uses NTT-friendly primes for which \(2N\)-th roots of
unity exist, so \(X^N+1\) may be reducible or fully split over the selected
prime field. This is operationally useful for fast NTTs, but it is also a
security-assumption choice. The Ring-LPN PCG paper explicitly distinguishes
irreducible and reducible/splittable instantiations, and later work points
out that programmable PCGs based on splittable Ring-LPN rely on newer
assumptions~\cite{boyle2022ringlpn,bombar2023qa}. Therefore a paper claim
must state which assumption is being used, the factorization degree of
\(X^N+1\) modulo each \(p_i\), and the concrete security estimate. It is not
scientifically sound to write simply ``secure under Ring-LPN'' without this
parameter mapping.

\section{The Share-Conversion Obstacle}

\RingLPN{} OLE naturally gives additive shares modulo an odd prime \(p\), or
per-limb shares that can be CRT-recomposed to a scalar modulo
\(M=p_0p_1\). \Orca{} FC layers consume additive shares in \(\modtwo\).
Naively reducing each share is incorrect.

Suppose \(z_0,z_1\in [0,M)\) and
\[
  v = z_0 + z_1 \pmod M.
\]
Let
\[
  m = \left\lfloor \frac{z_0+z_1}{M} \right\rfloor .
\]
The dealer/oracle conversion is
\[
\begin{aligned}
  r_0 &= z_0 \bmod 2^{\mathsf{bw}},\\
  r_1 &= z_1 - mM \bmod 2^{\mathsf{bw}}.
\end{aligned}
\]
Then
\[
  r_0+r_1 = v \pmod {2^{\mathsf{bw}}}.
\]
The hidden bit \(m\) is the barrier to dealerless conversion. A centralized
oracle can compute it because it sees both \(z_0\) and \(z_1\). In a true
dealerless protocol, neither party knows enough to compute \(m\) alone.

\paragraph{Required theory result.}
The project needs either:
\begin{enumerate}[leftmargin=*]
  \item a secure distributed wrap-bit/comparison protocol that outputs
  additive shares of \(mM \bmod 2^{\mathsf{bw}}\), or
  \item a different arithmetic design that keeps the online protocol in the
  prime/CRT ring until after reconstruction, or
  \item a cited PCG/OLE construction that directly emits correlations in
  \(\modtwo\).
\end{enumerate}
Until this boundary is solved, the phrase ``trusted dealer removed'' should
not be used.

This is a standard mixed-domain MPC issue rather than a coding detail:
modular conversion between additive shares over different moduli has dedicated
protocols~\cite{yu2011modcnv}, and modern mixed arithmetic/binary systems use
special preprocessing such as daBits/edaBits for bit-decomposition,
comparison, and conversion~\cite{escudero2020edabits}. Those works are not
drop-in implementations for this codebase, but they are the right literature
class for replacing the current oracle carry correction.

\section{Candidate Dealerless Protocol}

The following is the clean semi-honest protocol skeleton for one FC matmul.

\paragraph{Distributed setup.}
Both parties agree on \((N,\{p_i\},M,\mathsf{bw})\), a packing layout, and
layer dimensions \((M_{\mathrm{rows}},K,N_{\mathrm{cols}})\). They run a
distributed PCG/OLE setup protocol to obtain correlated seeds for the required
OLE instances. Each party also locally samples its mask shares:
\[
  \ServerZero: A_0,B_0,Y_0,\qquad
  \ServerOne: A_1,B_1,Y_1.
\]
If the setup is replaced by a local test oracle, the result is only a
simulation of the intended dealerless transcript.

\paragraph{Cross-term generation.}
For each ring product or packed product slot:
\[
\begin{aligned}
  (U_0,U_1) &\leftarrow \OLE(A_0,B_1),\\
  (V_0,V_1) &\leftarrow \OLE(A_1,B_0).
\end{aligned}
\]
The implementation should first use the existing Figure 2 SPFSS/OLE artifact
as the OLE engine, then replace any test oracle with the actual two-party
message boundary.

\paragraph{Local accumulation.}
Each party computes, per prime limb,
\[
\begin{aligned}
  Z_{0,i} &= A_{0,i}B_{0,i} + U_{0,i} + V_{0,i} \pmod {p_i},\\
  Z_{1,i} &= A_{1,i}B_{1,i} + U_{1,i} + V_{1,i} \pmod {p_i}.
\end{aligned}
\]
For q62 there is one limb. For q128, the limbs are later recomposed only for
the scalar export/carry reference.
Matrix products sum across the inner dimension. If constant-polynomial
packing is used, each scalar entry uses one polynomial product. If dense
packing is used, a separate extraction proof must identify the valid
coefficients and account for negacyclic sign flips.

\paragraph{Export to \Orca{}.}
Convert the per-limb \(Z_{0,i},Z_{1,i}\) shares into scalar shares modulo
\(M\), then into \(\modtwo\). In the current reference this is the
carry-corrected oracle conversion above. In the dealerless protocol this must
become a secure conversion. The parties then add their \(\modtwo\)
output-mask shares:
\[
  C_i = \mathsf{Convert}(Z_i) + Y_i \pmod {2^{\mathsf{bw}}}.
\]

\paragraph{Key writing.}
Each party writes its local key buffer:
\[
  \Party_i: \quad A_i \,\|\, B_i \,\|\, C_i.
\]
This must remain compatible with \texttt{FCLayer::readForwardKey} and
\texttt{readGPUMatmulKey}. The online path remains unchanged.

\section{Security Model}

The first written protocol should target semi-honest security. The proof
obligations are:
\begin{enumerate}[leftmargin=*]
  \item Local mask shares are sampled independently and reveal no information
  about the clear activation or weight values.
  \item The distributed PCG setup produces OLE seeds indistinguishable from
  ideal OLE correlations under the stated Ring-LPN or splittable Ring-LPN
  assumption.
  \item The OLE expansions reveal only additive shares of cross products.
  \item The conversion protocol from CRT/prime-limb shares to \(\modtwo\)
  shares leaks no information beyond the output shares.
  \item The final key buffers are distributed exactly as Beaver triples in
  \Orca{}'s existing offline dealer model.
\end{enumerate}

Malicious security would require additional consistency checks, especially
around OLE inputs and conversion. That should be treated as a later protocol
variant unless the professor asks for it explicitly.

\section{Implementation Roadmap}

\subsection*{Checkpoint 1: Freeze the Current Oracle Reference}

Preserve the current reference as an oracle specification:
\begin{itemize}[leftmargin=*]
  \item q62 and q128 host bridge tests must keep zero corrected failures;
  \item q62/full-32-bit negative control must remain present;
  \item the \Orca{} FC demo must keep passing q62 bounded and q128 bounded
  cases;
  \item feature flag off must preserve baseline \Orca{} behavior.
\end{itemize}

\subsection*{Checkpoint 2: Add an Ideal Dealerless Transcript Test}

Implement a host or CUDA-light test that samples \(A_0,A_1,B_0,B_1,Y_0,Y_1\)
separately, computes the two cross terms through an \emph{ideal OLE oracle},
and writes party-local key buffers. This is not the final protocol, but it is
the right executable specification:
\[
  \text{ideal OLE transcript} \Rightarrow
  \text{byte-compatible \Orca{} keys} \Rightarrow
  \texttt{gpuMatmulBeaver pass}.
\]
This test should report:
\[
  \#\OLE = 2\cdot \text{rows}\cdot K\cdot \text{cols}
\]
for constant-polynomial packing, and the number of conversion calls.

\subsection*{Checkpoint 3: Parameter and Assumption Audit}

Before replacing the oracle with the real Ring-LPN machinery, record the
cryptographic parameter interpretation:
\begin{itemize}[leftmargin=*]
  \item the exact primes \(p_i\), polynomial \(X^N+1\), and factorization
  behavior modulo each \(p_i\);
  \item whether the intended assumption is irreducible Ring-LPN, reducible
  Ring-LPN, or a splittable/divisible variant;
  \item the target security level and whether q128/actual q124 is sufficient
  for the target \Orca{} layer bounds;
  \item whether the protocol claim is semi-honest only or includes active
  consistency checks.
\end{itemize}

\subsection*{Checkpoint 4: Replace the Ideal OLE with Figure 2 OLE}

The next protocol-real step is to replace the ideal cross-term oracle with the
existing \RingLPN{} OLE engine:
\[
  \OLE(A_0,B_1),\qquad \OLE(A_1,B_0).
\]
Before claiming q128, regenerate saved q128 summaries for:
\begin{itemize}[leftmargin=*]
  \item Figure 2 OLE,
  \item ring-polynomial linear OLE-to-Beaver,
  \item \Orca{} FC keywriter with q128 bounds.
\end{itemize}

\subsection*{Checkpoint 5: Secure Share Conversion}

This is the central protocol gap. The deliverable should be a small standalone
document and a test implementation for secure conversion:
\[
  \Share{z}_{M} \longrightarrow \Share{z}_{2^{\mathsf{bw}}}.
\]
The oracle conversion remains the reference. A secure candidate must match it
bit-for-bit on randomized tests, forced-wrap tests, and layer-shaped matmul
outputs. Candidate approaches include secure wrap-bit computation, modular
conversion protocols, or mixed arithmetic/binary preprocessing such as
edaBits; the memo should explain why the chosen approach matches the two-party
\Orca{} threat model.

\subsection*{Checkpoint 6: Dense Packing}

Once correctness is stable, replace constant-polynomial packing. The packing
work needs its own proof:
\begin{enumerate}[leftmargin=*]
  \item define the coefficient layout for activation and weight tensors;
  \item prove which coefficients contain the desired dot products;
  \item handle negacyclic wraparound under \(X^N=-1\);
  \item validate packed and constant-polynomial outputs agree on small FC
  layers;
  \item report utilization as scalar products per \RingLPN{} OLE expansion.
\end{enumerate}

\subsection*{Checkpoint 7: Expand Beyond Forward FC}

After forward FC passes with the real protocol, extend to:
\[
  dW = X^\top dZ,\qquad dX = dZ W^\top.
\]
The current demo already has synthetic \(dW\) and \(dX\) contract checks, but
the full training path also needs truncation, optimizer, and bias-gradient
boundaries kept honest.

\section{Suggested Professor-Facing Claims}

Safe claims today:
\begin{enumerate}[leftmargin=*]
  \item The project has a byte-compatible \Orca{} FC keywriter demo.
  \item The unchanged \texttt{gpuMatmulBeaver} online path validates the
  generated key buffers.
  \item q128 CRT export is validated in the host/oracle bridge for bounded
  32-bit cases.
  \item The OLE-to-Beaver algebra is validated in standalone
  ring-polynomial matrix artifacts.
  \item The intended dealerless algebra is consistent with standard OLE-to-
  multiplication-triple reductions, provided the OLE and conversion
  functionalities are realized securely.
\end{enumerate}

Claims to avoid until additional work lands:
\begin{enumerate}[leftmargin=*]
  \item ``The trusted dealer is removed.''
  \item ``The conversion from \(\Z_M\) to \(\modtwo\) is secure.''
  \item ``Dense packing is implemented.''
  \item ``Full \Orca{} training uses Ring-LPN dealerless preprocessing.''
  \item ``Paper-parameter OLE/linear numbers are complete,'' unless the q128
  OLE and q128 linear summaries are regenerated and validated.
  \item ``Security follows from standard Ring-LPN'' without specifying
  whether the deployed NTT polynomial is irreducible, reducible, or fully
  split under the selected primes.
\end{enumerate}

\section{Concrete Next Deliverables}

\begin{table}[h]
\centering
\small
\begin{tabular}{p{0.08\linewidth}p{0.42\linewidth}p{0.40\linewidth}}
\toprule
Order & Deliverable & Evidence required \\
\midrule
1 & Ideal dealerless FC transcript & Key buffers pass unchanged
\texttt{gpuMatmulBeaver}; report OLE and conversion counts. \\
2 & Parameter/assumption audit & Prime limbs, factorization, assumption
variant, and layer bounds are written down. \\
3 & q128 OLE and q128 linear saved summaries & Uniform and regular smokes pass
with requested qbits 128, actual qbits 124. \\
4 & Secure conversion note and prototype & Matches oracle conversion on forced
wraps and layer-shaped randomized tests. \\
5 & Real OLE-backed FC transcript & Replaces ideal cross-term oracle with
Figure 2 OLE shares. \\
6 & Dense packing oracle & Packed and constant-polynomial paths agree on
small FC layers. \\
7 & Model-level FC benchmark & Feature flag on/off comparison for selected
\Orca{} FC-heavy workloads. \\
\bottomrule
\end{tabular}
\caption{Recommended order for the next dealerless work.}
\end{table}

\section{Summary}

The clean path is to keep \Orca{}'s online protocol fixed and make the offline
correlations progressively less oracle-like. The algebraic bridge is already
clear: two OLEs supply the cross terms that a dealer would otherwise compute.
The current codebase has enough machinery to validate byte compatibility and
bounded q128 export, but a true dealerless claim depends on replacing two
oracle boundaries: ideal cross-term generation and carry-corrected
\(\Z_M \rightarrow \modtwo\) conversion. It also depends on matching the
chosen NTT-friendly parameters to a clearly stated Ring-LPN assumption. The
next best checkpoint is therefore an ideal dealerless transcript test,
followed by a parameter audit, a secure-conversion design, and a real
\RingLPN{} OLE-backed transcript.

\begin{thebibliography}{9}

\bibitem{boyle2022ringlpn}
Elette Boyle, Geoffroy Couteau, Niv Gilboa, Yuval Ishai, Lisa Kohl, and
Peter Scholl.
\newblock Efficient Pseudorandom Correlation Generators from Ring-LPN.
\newblock IACR Cryptology ePrint Archive, Paper 2022/1035. Full version of
CRYPTO 2020 paper.
\newblock \url{https://eprint.iacr.org/2022/1035}.

\bibitem{boyle2019pcg}
Elette Boyle, Geoffroy Couteau, Niv Gilboa, Yuval Ishai, Lisa Kohl, and
Peter Scholl.
\newblock Efficient Pseudorandom Correlation Generators: Silent OT Extension
and More.
\newblock CRYPTO 2019.
\newblock \url{https://eprint.iacr.org/2019/448}.

\bibitem{baum2020ole}
Carsten Baum, Daniel Escudero, Alberto Pedrouzo-Ulloa, Peter Scholl, and
Juan Ramon Troncoso-Pastoriza.
\newblock Efficient Protocols for Oblivious Linear Function Evaluation from
Ring-LWE.
\newblock IACR Cryptology ePrint Archive, Paper 2020/970.
\newblock \url{https://eprint.iacr.org/2020/970}.

\bibitem{yu2011modcnv}
Ching-Hua Yu and Bo-Yin Yang.
\newblock Randomized Secure Two-Party Computation for Modular Conversion, Zero
Test, Comparison, MOD and Exponentiation.
\newblock IACR Cryptology ePrint Archive, Paper 2011/560.
\newblock \url{https://eprint.iacr.org/2011/560}.

\bibitem{escudero2020edabits}
Daniel Escudero, Satrajit Ghosh, Marcel Keller, Rahul Rachuri, and Peter
Scholl.
\newblock Improved Primitives for MPC over Mixed Arithmetic-Binary Circuits.
\newblock IACR Cryptology ePrint Archive, Paper 2020/338; CRYPTO 2020.
\newblock \url{https://eprint.iacr.org/2020/338}.

\bibitem{bombar2023qa}
Maxime Bombar, Geoffroy Couteau, Alain Couvreur, and Clement Ducros.
\newblock Correlated Pseudorandomness from the Hardness of Quasi-Abelian
Decoding.
\newblock arXiv:2306.03488, CRYPTO 2023 long version.
\newblock \url{https://arxiv.org/abs/2306.03488}.

\end{thebibliography}

\end{document}
